\documentclass[10pt,a4paper]{amsart}
\usepackage{amsmath,amssymb,amsthm,graphicx}
% Retained preamble for the Stacks Project chapter divisors.tex (Divisors).
% The chapter file pinned-original-source.tex carries only \input{preamble} on its first line, so no
% preamble text lives inside the chapter itself.  The Stacks Project uses one single shared preamble for
% every chapter; the snapshot stages it verbatim in sources/stacks-chow/preamble.tex.  The block below is
% copied mechanically, line for line, from that file's lines 14-18 (the Xy-pic setup) and lines 154-217
% (the native theorem declarations together with the mathematical macros).  Nothing was added, removed or
% redefined.  The wrapper class and the amsmath/amssymb/amsthm/graphicx base packages are supplied by the
% assembler, so the class/\externaldocument part of the shared preamble is not repeated here.
% For commutative diagrams we use Xy-pic
\usepackage[all]{xy}

% We use 2cell for 2-commutative diagrams.
\xyoption{2cell}
% Theorem environments.
%
\theoremstyle{plain}
\newtheorem{theorem}[subsection]{Theorem}
\newtheorem{proposition}[subsection]{Proposition}
\newtheorem{lemma}[subsection]{Lemma}

\theoremstyle{definition}
\newtheorem{definition}[subsection]{Definition}
\newtheorem{example}[subsection]{Example}
\newtheorem{exercise}[subsection]{Exercise}
\newtheorem{situation}[subsection]{Situation}

\theoremstyle{remark}
\newtheorem{remark}[subsection]{Remark}
\newtheorem{remarks}[subsection]{Remarks}

\numberwithin{equation}{subsection}

% Macros
%
\def\lim{\mathop{\mathrm{lim}}\nolimits}
\def\colim{\mathop{\mathrm{colim}}\nolimits}
\def\Spec{\mathop{\mathrm{Spec}}}
\def\Hom{\mathop{\mathrm{Hom}}\nolimits}
\def\Ext{\mathop{\mathrm{Ext}}\nolimits}
\def\SheafHom{\mathop{\mathcal{H}\!\mathit{om}}\nolimits}
\def\SheafExt{\mathop{\mathcal{E}\!\mathit{xt}}\nolimits}
\def\Sch{\mathit{Sch}}
\def\Mor{\mathop{\mathrm{Mor}}\nolimits}
\def\Ob{\mathop{\mathrm{Ob}}\nolimits}
\def\Sh{\mathop{\mathit{Sh}}\nolimits}
\def\NL{\mathop{N\!L}\nolimits}
\def\CH{\mathop{\mathrm{CH}}\nolimits}
\def\proetale{{pro\text{-}\acute{e}tale}}
\def\etale{{\acute{e}tale}}
\def\QCoh{\mathit{QCoh}}
\def\Ker{\mathop{\mathrm{Ker}}}
\def\Im{\mathop{\mathrm{Im}}}
\def\Coker{\mathop{\mathrm{Coker}}}
\def\Coim{\mathop{\mathrm{Coim}}}

% Boxtimes
%
\DeclareMathSymbol{\boxtimes}{\mathbin}{AMSa}{"02}

%
% Macros for moduli stacks/spaces
%
\def\QCohstack{\mathcal{QC}\!\mathit{oh}}
\def\Cohstack{\mathcal{C}\!\mathit{oh}}
\def\Spacesstack{\mathcal{S}\!\mathit{paces}}
\def\Quotfunctor{\mathrm{Quot}}
\def\Hilbfunctor{\mathrm{Hilb}}
\def\Curvesstack{\mathcal{C}\!\mathit{urves}}
\def\Polarizedstack{\mathcal{P}\!\mathit{olarized}}
\def\Complexesstack{\mathcal{C}\!\mathit{omplexes}}
% \Pic is the operator that assigns to X its picard group, usage \Pic(X)
% \Picardstack_{X/B} denotes the Picard stack of X over B
% \Picardfunctor_{X/B} denotes the Picard functor of X over B
\def\Pic{\mathop{\mathrm{Pic}}\nolimits}
\def\Picardstack{\mathcal{P}\!\mathit{ic}}
\def\Picardfunctor{\mathrm{Pic}}
\def\Deformationcategory{\mathcal{D}\!\mathit{ef}}

% --- appended 2026-10-12 (Chief): the project own definitions of the reference/slogan/history
% display environments, copied mechanically from the shared Stacks preamble lines 8-12
% (\usepackage{verbatim} plus the three \newenvironment lines, which the chapter preamble keeps
% verbatim from the project preamble). Nothing added, removed or redefined: these environments
% carry no mathematical content and the project itself typesets them as comments. Without these
% three lines a counted statement whose own text uses them fails with Environment slogan undefined.
\usepackage{verbatim}
\newenvironment{reference}{\comment}{\endcomment}
\newenvironment{slogan}{\comment}{\endcomment}
\newenvironment{history}{\comment}{\endcomment}
\usepackage{cleveref}
\crefname{originallemma}{original lemma}{original lemmas}
\Crefname{originallemma}{Original lemma}{Original lemmas}
\makeatletter
\newlabel{lemma-closed-subscheme-proj}{{0801}{source}{}{sourcepointer-1}{}}
\newlabel{lemma-closed-subscheme-proj@cref}{{[originallemma][0][]0801}{[1][0][]source}{}{}{}}
\makeatother
\crefname{originallemma}{original lemma}{original lemmas}
\Crefname{originallemma}{Original lemma}{Original lemmas}
\makeatletter
\newlabel{morphisms-lemma-closed-immersion-finite-presentation}{{01TV}{source}{}{sourcepointer-2}{}}
\newlabel{morphisms-lemma-closed-immersion-finite-presentation@cref}{{[originallemma][0][]01TV}{[1][0][]source}{}{}{}}
\makeatother
\crefname{originallemma}{original lemma}{original lemmas}
\Crefname{originallemma}{Original lemma}{Original lemmas}
\makeatletter
\newlabel{properties-lemma-quasi-coherent-colimit-finite-type}{{01PG}{source}{}{sourcepointer-3}{}}
\newlabel{properties-lemma-quasi-coherent-colimit-finite-type@cref}{{[originallemma][0][]01PG}{[1][0][]source}{}{}{}}
\makeatother
\title{A finitely presented closed subscheme of a relative Proj is cut out by finitely many forms}
\author{The Stacks Project Authors}
\date{stacks-project:divisors.tex:local-snapshot-20261012; GNU FDL 1.2 or later}
\begin{document}
\maketitle
\setcounter{subsection}{1}
\noindent\textit{Editorial note.} All delivered mathematical text is the author's own, copied from the pinned Stacks Project chapter file divisors.tex; no statement or proof step was rewritten, reordered or removed.\par\medskip
\noindent\textit{Editorial note.} Carrier window: statement and proof both lie inside the window 8000-10000 granted to this lane by CONTROL 8.83.\par\medskip
\noindent\textit{Editorial note.} The statement carries the underline-Proj notation of the chapter; it is part of the counted statement span and is preserved verbatim.\par\medskip
\noindent\textit{Editorial note.} The three references of the counted spans are recorded as non-load-bearing navigation pointers; the argument that consumes them stays inside the counted proof span.\par\medskip
\noindent\textit{Editorial note.} Preamble: sources/stacks-divisors/retained-preamble.tex (88 lines, verbatim present).\par\medskip
\noindent\textit{Allowed standard prerequisites (editorial).}\begin{itemize}
\item Closed subschemes of a relative Proj come from graded ideal sheaves. Lemma lemma-closed-subscheme-proj, used at proof lines 8110-8112. Supplies the graded ideal I with Z = Proj(A/I).
\item Finite presentation of a closed immersion means the ideal sheaf is of finite type. Morphisms, Lemma morphisms-lemma-closed-immersion-finite-presentation, used at proof lines 8123-8126. Produces the finitely many local generators h\_\{i,j,k\}.
\item Finite type submodules of a quasi-coherent module exist over quasi-compact bases. Properties, Lemma properties-lemma-quasi-coherent-colimit-finite-type, used at proof lines 8133-8136. Produces the finite type F inside I\_d.
\end{itemize}\medskip
\ifdefined\hypertarget\hypertarget{sourcepointer-1}{}\fi
\noindent\textit{Source reference.} Original lemma 0801: focused-seeds/AGP/sources/stacks-snapshot-20261011/tags.txt, source lines 10166--10166. Navigation pointer into the Stacks Project; the standard fact it supplies is applied inside the counted proof span and the argument resting on it stays there.\par\medskip
\ifdefined\hypertarget\hypertarget{sourcepointer-2}{}\fi
\noindent\textit{Source reference.} Original lemma 01TV: focused-seeds/AGP/sources/stacks-snapshot-20261011/tags.txt, source lines 2388--2388. Navigation pointer into the Stacks Project; the standard fact it supplies is applied inside the counted proof span and the argument resting on it stays there.\par\medskip
\ifdefined\hypertarget\hypertarget{sourcepointer-3}{}\fi
\noindent\textit{Source reference.} Original lemma 01PG: focused-seeds/AGP/sources/stacks-snapshot-20261011/tags.txt, source lines 2229--2229. Navigation pointer into the Stacks Project; the standard fact it supplies is applied inside the counted proof span and the argument resting on it stays there.\par\medskip
\begin{lemma}
\label{lemma-closed-subscheme-proj-finite}
Let $S$ be a quasi-compact and quasi-separated scheme.
Let $\mathcal{A}$ be a quasi-coherent graded $\mathcal{O}_S$-algebra. Let
$p : X = \underline{\text{Proj}}_S(\mathcal{A}) \to S$ be the relative
Proj of $\mathcal{A}$. Let $i : Z \to X$ be a closed subscheme.
If $p$ is quasi-compact and $i$ of finite presentation, then there exists
a $d > 0$ and a quasi-coherent finite type $\mathcal{O}_S$-submodule
$\mathcal{F} \subset \mathcal{A}_d$ such that
$Z = \underline{\text{Proj}}_S(\mathcal{A}/\mathcal{F}\mathcal{A})$.
\end{lemma}

\begin{proof}
By Lemma \ref{lemma-closed-subscheme-proj} we know there exists a
quasi-coherent graded sheaf of ideals $\mathcal{I} \subset \mathcal{A}$
such that $Z = \underline{\text{Proj}}(\mathcal{A}/\mathcal{I})$.
Since $S$ is quasi-compact we can choose a finite affine open covering
$S = U_1 \cup \ldots \cup U_n$. Say $U_i = \Spec(R_i)$. Let
$\mathcal{A}|_{U_i}$ correspond to the graded $R_i$-algebra $A_i$ and
$\mathcal{I}|_{U_i}$ to the graded ideal $I_i \subset A_i$. Note that
$p^{-1}(U_i) = \text{Proj}(A_i)$ as schemes over $R_i$.
Since $p$ is quasi-compact we can choose finitely many homogeneous
elements $f_{i, j} \in A_{i, +}$ such that
$p^{-1}(U_i) = \bigcup_j D_{+}(f_{i, j})$.
The condition on $Z \to X$ means that the ideal sheaf of $Z$ in
$\mathcal{O}_X$ is of finite type, see
Morphisms, Lemma \ref{morphisms-lemma-closed-immersion-finite-presentation}.
Hence we can find finitely many homogeneous elements
$h_{i, j, k} \in I_i \cap A_{i, +}$ such that the ideal of
$Z \cap D_{+}(f_{i, j})$ is generated by the elements
$h_{i, j, k}/f_{i, j}^{e_{i, j, k}}$. Choose $d > 0$ to be a common multiple
of all the integers $\deg(f_{i, j})$ and $\deg(h_{i, j, k})$.
By Properties, Lemma \ref{properties-lemma-quasi-coherent-colimit-finite-type}
there exists a finite type quasi-coherent $\mathcal{F} \subset \mathcal{I}_d$
such that all the local sections
$$
h_{i, j, k}f_{i, j}^{(d - \deg(h_{i, j, k}))/\deg(f_{i, j})}
$$
are sections of $\mathcal{F}$. By construction $\mathcal{F}$ is a solution.
\end{proof}

\end{document}
